\documentclass[11pt,a4paper]{article}
\usepackage[T1]{fontenc}
\usepackage[utf8]{inputenc}
\usepackage{lmodern,amsmath,amssymb}
\usepackage[margin=25mm]{geometry}
\usepackage[hidelinks]{hyperref}
\hypersetup{pdftitle={The Planck gap as a theorem: every marking protocol needs $F^2\tau^3>9m\kappa$},pdfauthor={navstokgap research working notes}}
\newcommand{\tightlist}{\setlength{\itemsep}{0pt}\setlength{\parskip}{0pt}}
\setlength{\emergencystretch}{3em}
\setcounter{secnumdepth}{0}
\begin{document}
\hypertarget{the-planck-gap-as-a-theorem-every-marking-protocol-needs-f2tau39mkappa}{%
\section{\texorpdfstring{The Planck gap as a theorem: every marking
protocol needs
\(F^2\tau^3>9m\kappa\)}{The Planck gap as a theorem: every marking protocol needs F\^{}2\textbackslash tau\^{}3\textgreater9m\textbackslash kappa}}\label{the-planck-gap-as-a-theorem-every-marking-protocol-needs-f2tau39mkappa}}

Let \(\kappa\) be the smallest product of resolution and recoil width
over the marks a record model can make. Then for \textbf{every} finite
protocol of marks, of any number, at any times, the falling parabola can
be told from the inertial line over a duration \(\tau\) in the worst
case only if

\[\boxed{F^2\tau^3>9\,m\kappa,\qquad\text{equivalently}\qquad
\tau\Delta E=\frac{F^2\tau^3}{2m}>\frac92\,\kappa,\qquad
A_{\rm inertial,fall}=\frac{vF\tau^3}{6m}>\frac{3v\kappa}{F},}\]

and a two-mark protocol decides as soon as \(F^2\tau^3>36m\kappa\) when
marks saturating \(\kappa\) are available (Theorems 1--2). The threshold
duration is therefore pinned within a factor \(4^{1/3}\) by \(\kappa\)
alone, and the Planck gap of the recorded Galileo comparison is positive
exactly when \(\kappa>0\) (Corollary 3). The same bound governs
Democritus insertion: the marks inside a window of duration \(\tau'\),
used by themselves, record the force only if \(F^2\tau'^3>9m\kappa\)
(Corollary 4). Any record model in which a probe's later position
records fix its momentum to arbitrary accuracy has \(\kappa=0\)
(Proposition 5), and Newton's optics gives the heuristic value
\(\kappa=\Lambda p=h/2\) of the \href{newton-mark-floor.md}{mark-floor
note}. \textbf{Quantum kinematics has no non-vacuous instance of this
framework:} no mark has both a finite worst-case resolution and a finite
worst-case recoil width, by Paley--Wiener, so \(\kappa=\infty\) and the
theorem holds emptily (Proposition 6). The result is therefore a theorem
about classical bounded-error record models, sharpening the C068--C123
closures into a characterization, and the bridge to \(\hbar\) requires
the probabilistic reformulation of Section 6, where the inequality
becomes the standard quantum limit for free-mass force detection and
must be argued against the Yuen--Caves--Ozawa dispute over that limit.
The proof is a separation argument for the adversary's feasibility set,
an AM--GM step that turns \(\delta\Delta\ge\kappa\) into a square root,
and one integration by parts that converts the certificate into a
Dirichlet energy. Constants are explicit throughout; nothing is fitted.

\hypertarget{setting-and-definitions}{%
\subsection{1. Setting and definitions}\label{setting-and-definitions}}

Transverse coordinate \(y\), mass \(m>0\), comparison on \([0,\tau]\)
between the inertial hypothesis \(\mathrm I\) and the falling hypothesis
\(\mathrm F\), which differ by the displacement

\[P(t)=\frac{Ft^2}{2m},\qquad P(0)=P'(0)=0,\qquad P''=\frac Fm .\]

A \textbf{mark} at time \(t_j\) has a \textbf{resolution}
\(\delta_j>0\): the datum it leaves is an interval of length
\(\delta_j\) containing \(y(t_j)\); and a \textbf{recoil width}
\(\Delta_j\ge0\): the transverse impulse it delivers to the body lies,
under either hypothesis, in an interval of length \(\Delta_j\) that is
fixed by the totality of records, so that two runs which agree in every
record can differ in that impulse by at most \(\Delta_j\). Between marks
the body moves freely under \(\mathrm I\) and with the force \(F\) under
\(\mathrm F\). A \textbf{protocol} is a finite set of marks at times
\(0=t_0<t_1<\dots<t_k\le\tau\); the mark at \(t_0=0\) is the
preparation, whose recoil width \(\Delta_0\) is the width within which
the records fix the initial transverse momentum. A protocol
\textbf{decides} the comparison in the worst case when no pair of runs,
one under \(\mathrm I\) and one under \(\mathrm F\), with initial
positions, initial momenta and recoils allowed by the marks, produces
the same data. Set

\[\kappa=\inf_j\delta_j\Delta_j\]

over the marks the model can make; a model satisfies the \textbf{mark
trade-off} with constant \(\kappa\) when every available mark has
\(\delta\Delta\ge\kappa\). For the corpuscular model of the mark-floor
note this is premise M3 with \(\kappa=\Lambda p\).

\textbf{Feasibility system.} Write \(y_{\rm F}-y_{\rm I}=P+d\). Then
\(d(t)=x+vt+\sum_{j\ge1}\rho_j(t-t_j)_+/m\) with \(x\in\mathbb R\) free
(both initial positions lie in the same datum interval, so their
difference is bounded only by the datum constraint at \(t_0\)),
\(|v|\le\Delta_0/m\) and \(|\rho_j|\le\Delta_j\). Two runs give the same
data exactly when

\[|d(t_i)+P(t_i)|\le\delta_i\qquad(i=0,1,\dots,k). \tag{1}\]

The protocol decides if and only if (1) has no solution
\((x,v,\rho_1,\dots,\rho_{k-1})\) in the stated ranges. (The last mark's
recoil affects nothing that is recorded and is omitted.)

\hypertarget{certificate-of-decidability}{%
\subsection{2. Certificate of
decidability}\label{certificate-of-decidability}}

\textbf{Lemma 1 (separation).} The protocol decides if and only if there
are multipliers \(\mu_0,\dots,\mu_k\in\mathbb R\) with \(\sum_i\mu_i=0\)
and

\[\sum_i\mu_iP(t_i)>\sum_i|\mu_i|\delta_i+\frac{\Delta_0}{m}|S_0|
+\sum_{j\ge1}\frac{\Delta_j}{m}|S_j|,\qquad
S_j=\sum_{i>j}\mu_i(t_i-t_j). \tag{2}\]

\emph{Proof.} Let \(L:(x,v,\rho)\mapsto(d(t_i))_i\in\mathbb R^{k+1}\) be
the linear map of Section 1 and \(B\) the box \(|v|\le\Delta_0/m\),
\(|\rho_j|\le\Delta_j\). The system (1) is feasible exactly when the
closed convex set \(C=\{L(x,v,\rho):x\in\mathbb R,(v,\rho)\in B\}\)
meets the compact box \(Q=\{q:|q_i+P(t_i)|\le\delta_i\}\). \(C\) is the
sum of a line and a compact convex set, hence closed and convex. If
\(C\cap Q=\emptyset\), the separation theorem for a closed and a compact
convex set gives \(\mu\) with
\(\sup_{q\in Q}\mu\cdot q<\inf_{c\in C}\mu\cdot c\). Since \(x\) is free
the infimum is \(-\infty\) unless \(\mu\cdot L(1,0,0)=\sum_i\mu_i=0\);
with that,
\(\inf_C\mu\cdot c=-\frac{\Delta_0}m|\sum_i\mu_it_i|-\sum_{j\ge1}\frac{\Delta_j}m|\sum_{i>j}\mu_i(t_i-t_j)|\)
by choosing the signs of \(v\) and \(\rho_j\), and
\(\sup_Q\mu\cdot q=-\sum_i\mu_iP(t_i)+\sum_i|\mu_i|\delta_i\). The
strict inequality between them is (2); if the separation comes with the
opposite orientation, replace \(\mu\) by \(-\mu\). Note
\(\sum_i\mu_it_i=S_0\) when \(\sum_i\mu_i=0\). Conversely, if (2) holds
and \((x,v,\rho)\) solved (1), pairing (1) with \(\mu\) would give
\(\sum_i\mu_iP(t_i)\le\sum_i|\mu_i|\delta_i-\sum_i\mu_id(t_i) \le\sum_i|\mu_i|\delta_i+\frac{\Delta_0}m|S_0|+\sum_{j\ge1}\frac{\Delta_j}m|S_j|\),
contradicting (2). \(\square\)

\textbf{Lemma 2 (trade-off to a square root).} If
\(\delta_j\Delta_j\ge\kappa\) for all \(j\), then the right side of (2)
is at least \(2\sqrt{\kappa/m}\,\sum_{j=0}^k\sqrt{|\mu_j||S_j|}\).

\emph{Proof.} For each \(j\),
\(|\mu_j|\delta_j+\frac{\Delta_j}m|S_j| \ge2\sqrt{|\mu_j||S_j|\delta_j\Delta_j/m}\ge2\sqrt{\kappa|\mu_j||S_j|/m}\)
by the arithmetic--geometric mean inequality; \(S_k=0\). \(\square\)

\hypertarget{the-universal-lower-bound}{%
\subsection{3. The universal lower
bound}\label{the-universal-lower-bound}}

\textbf{Lemma 3 (the certificate as a Dirichlet energy).} Let
\(\sum_i\mu_i=0\), \(N(u)=\sum_{i:\,t_i>u}\mu_i\) and
\(T(u)=\int_u^\tau N(w)\,dw\) for \(u\in\mathbb R\). Then \(T\) is
continuous and piecewise linear, \(T\equiv T(0)\) on \((-\infty,0]\),
\(T\equiv0\) on \([\tau,\infty)\), \(T(t_j)=S_j\), the slope of \(T\)
jumps by \(\mu_j\) at \(t_j\), and

\[\sum_i\mu_iP(t_i)=\frac Fm\int_0^\tau T(u)\,du,\qquad
\sum_j\mu_jT(t_j)=-\int_0^\tau T'(u)^2\,du. \tag{3}\]

\emph{Proof.} \(N\) is piecewise constant, vanishes for \(u<0\) because
\(\sum_i\mu_i=0\) and for \(u\ge t_k\), and drops by \(\mu_j\) as \(u\)
crosses \(t_j\); so \(T'=-N\) has the stated jumps and support. Also
\(T(t_j)=\sum_i\mu_i\,|\{u\in[t_j,\tau]:u<t_i\}|=\sum_{i>j}\mu_i(t_i-t_j)=S_j\).
For the first identity, \(P(t_i)=\int_0^{t_i}P'(u)\,du\) and Fubini give
\(\sum_i\mu_iP(t_i)=\int_0^\tau P'(u)N(u)\,du=-\int_0^\tau P'T'\,du =-[P'T]_0^\tau+\int_0^\tau P''T\,du=\frac Fm\int_0^\tau T\,du\),
using \(T(\tau)=0\) and \(P'(0)=0\). For the second, \(T'\) is of
bounded variation with jumps \(\mu_j\) at \(t_j\) and no other
variation, so
\(\sum_j\mu_jT(t_j)=\int_{\mathbb R}T\,dT'=[TT']_{-\infty}^{\infty}-\int T'^2 =-\int_0^\tau T'^2\),
as \(T'\) vanishes outside \([0,\tau]\). \(\square\)

\textbf{Theorem 1 (universal lower bound).} Suppose every mark of a
protocol satisfies \(\delta_j\Delta_j\ge\kappa\). If the protocol
decides the Galileo comparison on \([0,\tau]\) in the worst case, then

\[F^2\tau^3>9\,m\kappa .\]

\emph{Proof.} By Lemma 1 there is \(\mu\) with \(\sum_i\mu_i=0\)
satisfying (2), and by Lemma 2 and (3)

\[\frac Fm\int_0^\tau T\,du>2\sqrt{\frac\kappa m}\sum_j\sqrt{|\mu_j||S_j|}
\ge2\sqrt{\frac\kappa m}\sqrt{\sum_j|\mu_j||S_j|}
\ge2\sqrt{\frac\kappa m}\sqrt{E},\qquad E=\int_0^\tau T'^2\,du,\]

using \(\sum_j\sqrt{a_j}\ge\sqrt{\sum_ja_j}\) for \(a_j\ge0\) and
\(\sum_j|\mu_j||S_j|\ge|\sum_j\mu_jT(t_j)|=E\). The right side is
nonnegative, so \(\int_0^\tau T>0\). Since \(T(\tau)=0\),
\(|T(u)|=|\int_u^\tau T'|\le\sqrt{\tau-u}\sqrt E\) by Cauchy--Schwarz,
hence
\(\int_0^\tau T\le\int_0^\tau\sqrt{\tau-u}\,du\,\sqrt E=\tfrac23\tau^{3/2}\sqrt E\),
that is \(\sqrt E\ge\tfrac32\tau^{-3/2}\int_0^\tau T\). Inserting,
\(\frac Fm\int_0^\tau T>3\sqrt{\kappa/m}\,\tau^{-3/2}\int_0^\tau T\),
and dividing by the positive integral, \(F\tau^{3/2}>3\sqrt{\kappa m}\).
\(\square\)

The theorem uses nothing about the marks except the product bound: no
least length, no least impulse separately, no restriction on how many
marks are made or when, no assumption that the marks are of one kind,
and no probabilistic model. In the mark-floor note's two-mark protocol
the same inequality appeared with the constant \(16\); Theorem 1 says
that no protocol whatever can push the constant below \(9\).

\hypertarget{the-matching-upper-bound-and-the-gap}{%
\subsection{4. The matching upper bound and the
gap}\label{the-matching-upper-bound-and-the-gap}}

\textbf{Theorem 2 (two marks suffice).} Suppose the model offers, for
the duration \(\tau\) in question, a preparation mark with
\(\delta_0=\sqrt{\kappa\tau/m}\) and \(\Delta_0=\kappa/\delta_0\), and a
terminal mark with \(\delta_1\le\sqrt{\kappa\tau/m}\). Then the two-mark
protocol decides the comparison whenever \(F^2\tau^3>36\,m\kappa\).

\emph{Proof.} Under \(\mathrm I\) the terminal position lies in an
interval of length \(w=\delta_0+\Delta_0\tau/m=2\sqrt{\kappa\tau/m}\);
under \(\mathrm F\) in its translate by \(s(\tau)=F\tau^2/2m\). A datum
of length \(\delta_1\) cannot meet both when \(s(\tau)>w+\delta_1\),
which holds if \(F\tau^2/2m>3\sqrt{\kappa\tau/m}\), that is
\(F^2\tau^3>36m\kappa\). \(\square\)

\textbf{Corollary 3 (the Planck gap of the recorded comparison).} Let
\(\tau_*(F,m)\) be the infimum of durations over which some available
protocol decides the comparison. Under the hypotheses of Theorems 1--2,

\[9\,m\kappa\le F^2\tau_*^3\le36\,m\kappa,\qquad
\frac92\kappa\le\tau_*\Delta E(\tau_*)\le18\kappa,\qquad
\frac{3v\kappa}{F}\le A(\tau_*)\le\frac{12v\kappa}{F},\]

the lower bounds from Theorem 1 as infima and the upper bounds from
Theorem 2.

In particular the recorded comparison has a positive floor if and only
if \(\kappa>0\), and the floor is a fixed multiple of \(\kappa\),
between \(\tfrac92\) and \(18\) in the action-scaled variable
\(\tau\Delta E\). The geometric area inherits the factor \(v/F\) and has
no floor of its own, which is the
\href{action-unit-dimensional-selection.md}{Q14} gate: the only action
the record model contains is \(\kappa\).

\textbf{Corollary 4 (Democritus insertion).} Let marks be inserted at
\(t_a=t_{j}<t_{j+1}<\dots<t_{j+r}=t_b\) inside \([0,\tau]\), and let the
observer use these marks alone to decide whether the motion on
\([t_a,t_b]\) was free or forced, with no information about the momentum
at \(t_a\) beyond what these marks give. Then the window records the
force only if \(F^2(t_b-t_a)^3>9m\kappa\).

\emph{Proof.} Shift the origin to \(t_a\). With the momentum at \(t_a\)
unknown, \(v\) is free as well as \(x\), so the certificate of Lemma 1
acquires the extra condition \(S_0=\sum_i\mu_it_i=0\) and loses the term
\(\frac{\Delta_0}m|S_0|\), which was zero under that condition anyway.
Lemmas 2--3 and the proof of Theorem 1 go through unchanged with
\(\tau\) replaced by \(t_b-t_a\). \(\square\)

So inserted points finer than \(\tau_*=(9m\kappa/F^2)^{1/3}\) show, by
themselves, free motion only; the force reappears when enough of them
are pooled that the pooled window exceeds \(\tau_*\), and it also
reappears if momentum information from before the window is carried in,
in which case the comparison is the whole-interval one and Theorem 1
applies to it. This is the rigorous form of the mark-floor note's
Theorem 2, and of I003's remark 3: \(\kappa\) controls the mesh at which
the marked polygon stops converging to Newton's curve.

\hypertarget{what-fixes-kappa}{%
\subsection{\texorpdfstring{5. What fixes
\(\kappa\)}{5. What fixes \textbackslash kappa}}\label{what-fixes-kappa}}

\textbf{Proposition 5 (deterministic probes give \(\kappa=0\)).} Suppose
the marks are made by probes that travel freely after the mark, that a
probe's transverse momentum after the mark equals its momentum before
minus the impulse delivered, and that a probe's position can be recorded
at any distance \(D\) with a resolution \(\delta_s\) bounded
independently of \(D\). Then for every mark
\(\Delta\le p_\parallel(\delta+\delta_s)/D\) for all \(D\), where
\(p_\parallel\) is the probe's longitudinal momentum, so \(\Delta=0\)
and \(\kappa=0\).

\emph{Proof.} The probe's transverse momentum after the mark is
\(p_\parallel\) times the tangent of its direction, which the two
recorded positions fix to within \((\delta+\delta_s)/D\); momentum
conservation transfers the same width to the delivered impulse; let
\(D\to\infty\). \(\square\)

This is the far-screen countermodel of the mark-floor note stated for
any deterministic probe, and it is the mechanism behind the classical
closures C068--C123 (\href{action-scale-obstructions.md}{synthesis}):
full-state records fix the recoil, the trade-off constant is zero, and
Corollary 3 gives floor zero. A positive \(\kappa\) therefore requires
that the probe's later position records \textbf{do not} fix its
momentum, which is an indeterminacy of the probe's own kinematics.

\textbf{Proposition 6 (the worst-case framework has no quantum
instance).} In quantum kinematics no mark has both a finite worst-case
resolution and a finite worst-case recoil width. Hence
\(\kappa=\infty\), Theorem 1 holds emptily, and no protocol of marks
decides the comparison with worst-case certainty.

\emph{Proof.} Suppose a mark's datum confines the position to an
interval \(I\) of length \(\delta<\infty\) with certainty, so that every
conditional post-mark state \(\psi\) has
\(\operatorname{supp}\psi\subseteq I\). By the Paley--Wiener theorem
\(\hat\psi\) extends to an entire function of exponential type; a
nonzero entire function cannot vanish on a set with nonempty interior,
so the momentum distribution \(|\hat\psi|^2\) has full support. The
impulse the mark delivered is the difference between the post-mark and
pre-mark momenta, so no bounded interval fixed by the records contains
it with certainty, and \(\Delta=\infty\). Symmetrically, a mark with
\(\Delta<\infty\) confines momentum to a bounded set with certainty and
so has \(\delta=\infty\). \(\square\)

This is the decisive limitation of the present theorem, and it is
structural rather than technical. Worst-case interval widths express a
record model in which a datum \emph{excludes} values outright. Quantum
data exclude nothing: every position outcome leaves momentum tails, and
every finite sequence of finite-precision quantum measurements leaves
the two hypotheses compatible with the data at some probability. So
Theorem 1, Theorem 2 and Corollaries 3--4 are theorems about
\textbf{classical bounded-error record models} --- the family of
C068--C123, whose closures they sharpen into a characterization, with
Proposition 5 supplying the mechanism. The claim in an earlier version
of this note that the quantum comparison needs \(\tau\Delta E>9\hbar\)
is withdrawn: it read \(\kappa\ge2\hbar\) off an error--disturbance
relation whose quantities are not the worst-case widths used here, and
the correct worst-case value is \(\infty\).

\textbf{What the probabilistic version must be, and against what.}
Replace the decision criterion by a hypothesis test: two prepared
ensembles, a fixed error probability \(\epsilon\), and a criterion that
the data distinguish \(\mathrm I\) from \(\mathrm F\) at level
\(\epsilon\). Replace \(\delta\) and \(\Delta\) by the error and
disturbance of an instrument in a stated metric. The expected shape of
the conclusion, \(F^2\tau^3\gtrsim m\hbar\), is the standard quantum
limit for detecting a force on a free mass
(\href{https://doi.org/10.1017/CBO9780511622748}{Braginsky and Khalili,
\emph{Quantum Measurement}, 1992}, metadata). That limit is contested in
exactly the regime this note needs: Yuen argued it can be beaten,
\href{https://doi.org/10.1103/PhysRevLett.54.2465}{Caves defended it}
(PRL \textbf{54}, 2465, 1985, abstract), and
\href{https://doi.org/10.1103/PhysRevLett.60.385}{Ozawa exhibited a
measurement breaking the standard quantum limit for free-mass position}
(PRL \textbf{60}, 385, 1988, abstract) using an instrument whose
disturbance is correlated with its outcome rather than independent of
it. The same distinction reappears in the error--disturbance literature,
where Ozawa's counterexamples to the naive Heisenberg product coexist
with the proved calibration-based relation of
\href{https://doi.org/10.1103/PhysRevLett.111.160405}{Busch, Lahti and
Werner} (PRL \textbf{111}, 160405, 2013, abstract;
\href{https://doi.org/10.1063/1.4871444}{J. Math. Phys. \textbf{55},
042111, 2014}, abstract). A probabilistic version of Theorem 1 is
therefore not a corollary of any single error--disturbance inequality.
It has to say which instruments are admitted, and its interest lies
precisely in whether the protocol-universality proved here (over all
finite numbers of marks at all times) survives when Ozawa's
outcome-correlated disturbances are admitted.

\textbf{The Newton-age value.} The corpuscular premises M1--M3 of the
mark-floor note give \(\kappa=\Lambda p\), with \(\Lambda=1/89000\) inch
measured by Newton and \(p\) unmeasured in his age; on the
identification \(p=h/\lambda\), \(\Lambda=\lambda/2\) this is
\(\kappa=h/2\). Newton-age materials give Corollary 3 in full once M3 is
granted, and Proposition 5 shows that Newton's own determinate fits deny
M3 and give \(\kappa=0\). The corpuscular model is a classical
bounded-error record model, which is why it fits this framework and why
its value of \(\kappa\) is a heuristic analogy rather than a quantum
result. The logical structure established here is therefore:

\[\text{floor of the classically recorded comparison}\;>0
\quad\Longleftrightarrow\quad\kappa>0,\qquad
\text{deterministic probes}\;\Longrightarrow\;\kappa=0,\]

proved here for all finite mark protocols with explicit constants. No
step assumes a path-integral phase rule, a fixed measurement budget, or
a supplied action unit other than \(\kappa\) itself, and \(\kappa\) is
characterized rather than supplied on the classical side. The quantum
side is open: Proposition 6 shows this framework cannot reach it, and
names the reformulation that could.

\hypertarget{consequence-for-state}{%
\subsection{6. Consequence for STATE}\label{consequence-for-state}}

What is proved: Theorems 1--2 sandwich the threshold of the recorded
Galileo comparison between \(9m\kappa\) and \(36m\kappa\) in
\(F^2\tau^3\) for every finite protocol of marks with
\(\delta\Delta\ge\kappa\); Corollary 4 does the same for a window of
inserted marks used by itself; Proposition 5 gives \(\kappa=0\) for
every deterministic probe with far-field position records. Together
these characterize the floor of a classical bounded-error record model
by one number and identify what kills it.

What is not proved, and was wrongly claimed in the first version of this
note: any quantum value of \(\kappa\). Proposition 6 shows the
worst-case framework has no non-vacuous quantum instance, so the route
to \(\hbar\) runs through a probabilistic reformulation, which lands in
the disputed territory of the standard quantum limit for free-mass
measurement (Yuen; Caves 1985; Ozawa 1988) and of error--disturbance
relations (Ozawa; Busch--Lahti--Werner 2013).

Open, in the order that matters for a publishable result:

\begin{enumerate}
\def\labelenumi{\arabic{enumi}.}
\tightlist
\item
  \textbf{The probabilistic theorem.} Restate Theorem 1 with a
  hypothesis test at fixed error probability and instrument-theoretic
  error and disturbance. Decide whether protocol-universality survives
  outcome-correlated disturbances of Ozawa's kind. This is the step that
  would turn the note into a foundations paper.
\item
  \textbf{The sharp constant.} The interval \([9,36]\) in
  \(F^2\tau_*^3/(m\kappa)\) is not closed; find the optimal protocol and
  the matching certificate.
\item
  \textbf{General force law.} Extend from \(P=Ft^2/2m\) to general \(P\)
  with \(P(0)=P'(0)=0\); the proof already pairs \(\int_0^\tau P''T\)
  against the Dirichlet energy of \(T\), so the natural statement bounds
  a norm of \(P''\) from below.
\item
  \textbf{Why \(\kappa>0\) without quantum kinematics.} Proposition 5
  now names exactly what must fail: far-field position records fixing a
  probe's momentum.
\item
  \textbf{Prior art.} Position the inequality explicitly as a
  probability-free, protocol-universal relative of the standard quantum
  limit rather than as a new bound.
\end{enumerate}
\end{document}
